% ============================================================================
% CHAPITRE 3 - LE SEMANTIC ONE-TIME PAD (bloc issu de clo.tex)
% ============================================================================

\section{LE SEMANTIC ONE-TIME PAD~: CONSTRUCTION DE BASE}
\label{sec:sotp_base}

% Bloc issu de clo.tex : paradigme Gold-Iron-Clay, S-OTP de base, metriques CLO atteintes
\subsection{Le paradigme Gold-Iron-Clay}
\label{subsec:gold_iron_clay}
% -----------------------------------------------------------------------
\begin{sloppypar}
Les trois m\'etriques \'etant d\'efinies, il faut une m\'ethodologie pour
concevoir des protocoles qui les \'equilibrent~\cite{Shoup01,Damgard87}. Le paradigme
\textbf{Gold-Iron-Clay} structure cette conception en trois couches
fonctionnelles. Sa d\'enomination, inspir\'ee de la m\'etallurgie m\'edi\'evale,
refl\`ete la relation entre mall\'eabilit\'e et r\'esistance.

\paragraph{Couche Gold (confidentialit\'e).}
Masquage cryptographique du message via one-time pads bas\'es sur PRG,
chiffrements s\'emantiquement s\^urs~\cite{GoldwasserMicali84,Bellare97}, ou preuves zero-knowledge~\cite{GMR85,BenOr88}. La couche
Gold r\'epond \`a la question \emph{\guillemotleft~Comment rendre $\sigma$ indiscernable entre deux
messages~?\guillemotright}

\paragraph{Couche Iron (fiabilit\'e).}
Commitments contraignants (\emph{binding}) permettant la r\'ecup\'eration
v\'erifiable, commitments de Pedersen~\cite{Pedersen91,ChaumPedersen92}, accumulateurs~\cite{CamenischLysyanskaya02}, signatures. La couche
Iron r\'epond \`a la question \emph{\guillemotleft~Comment garantir que $\sigma$
permet de retrouver $m$ avec probabilit\'e $1 - \negl(\lambda)$~?\guillemotright}

\paragraph{Couche Clay (opposabilit\'e).}
Encodage pr\'eservant le contexte institutionnel pour l'extraction
d'information par des interpr\`etes aux ressources born\'ees~\cite{Miller56,Simon82}, par encodage
d\'eterministe, compression Johnson-Lindenstrauss~\cite{JohnsonLindenstrauss84,DasguptaGupta03} ou hachage d\'eterministe. La couche
Clay r\'epond \`a la question \emph{\guillemotleft~Comment permettre \`a $\mathsf{Interpret}$ d'extraire
un maximum d'information sur $V$ depuis $(c, \sigma)$~?\guillemotright}

\begin{figure}[htbp]
\centering
\begin{tikzpicture}[
    scale=0.9,
    transform shape
]

% ---------------------------------------------------------------
% D\'EGRAD\'ES M\'ETALLIQUES
% ---------------------------------------------------------------

% Couche Or
\shade[left color=metalorfonce!60, right color=metalor!40,
       middle color=metalor!80]
    (-4, 4) rectangle (4, 5.5);
\draw[thick, metalorfonce] (-4, 4) rectangle (4, 5.5);

% Couche Fer
\shade[left color=metalferfonce!70, right color=metalfer!50,
       middle color=metalfer!90]
    (-4, 2) rectangle (4, 3.5);
\draw[thick, metalferfonce] (-4, 2) rectangle (4, 3.5);

% Couche Argile
\shade[left color=metalargilefonce!70, right color=metalargile!50,
       middle color=metalargile!90]
    (-4, 0) rectangle (4, 1.5);
\draw[thick, metalargilefonce] (-4, 0) rectangle (4, 1.5);

% ---------------------------------------------------------------
% \'ETIQUETTES DES COUCHES
% ---------------------------------------------------------------
\node[font=\bfseries\sffamily, text=metalorfonce] at (0, 4.75)
    {Couche Or \quad - \quad Confidentialit\'e};
\node[font=\bfseries\sffamily, text=white] at (0, 2.75)
    {Couche Fer \quad - \quad Fiabilit\'e};
\node[font=\bfseries\sffamily, text=white] at (0, 0.75)
    {Couche Argile \quad - \quad Opposabilit\'e};

% ---------------------------------------------------------------
% CONTENU DES COUCHES
% ---------------------------------------------------------------
\node[font=\small, align=center, text=metalorfonce] at (0, 4.25)
    {Masquage PRG, preuves ZK \quad $\Conf \to 1$};
\node[font=\small, align=center, text=white] at (0, 2.25)
    {Engagement liant \quad $\Relver \to 1$};
\node[font=\small, align=center, text=white] at (0, 0.25)
    {Encodage contextuel \quad $\Opp \to \max$};

% ---------------------------------------------------------------
% REFLETS M\'ETALLIQUES
% ---------------------------------------------------------------
\draw[white!70!metalor, very thick, opacity=0.4]
    (-3.6, 5.2) -- (3.6, 5.2);
\draw[white!50!metalfer, very thick, opacity=0.3]
    (-3.6, 3.2) -- (3.6, 3.2);
\draw[white!50!metalargile, very thick, opacity=0.3]
    (-3.6, 1.2) -- (3.6, 1.2);

% ---------------------------------------------------------------
% FL\`ECHES DE FLUX
% ---------------------------------------------------------------
\draw[-{Stealth[length=2.5mm]}, thick, metalorfonce]
    (-5.2, 4.75) -- (-4.2, 4.75)
    node[left, xshift=-1cm, font=\small] {$M$};
\draw[-{Stealth[length=2.5mm]}, thick, metalferfonce]
    (4.2, 2.75) -- (5.2, 2.75)
    node[right, font=\small] {$\Sigma$};
\draw[-{Stealth[length=2.5mm]}, thick, metalargilefonce]
    (4.2, 0.75) -- (5.2, 0.75)
    node[right, font=\small] {$V$};

% ---------------------------------------------------------------
% FL\`ECHES VERTICALES INTER-COUCHES
% ---------------------------------------------------------------
\draw[-{Stealth[length=2.5mm]}, thick, chaptercolor]
    (3.2, 3.9) -- (3.2, 3.6);
\draw[-{Stealth[length=2.5mm]}, thick, chaptercolor]
    (3.2, 1.9) -- (3.2, 1.6);

% ---------------------------------------------------------------
% ACCOLADES LAT\'ERALES (propri\'et\'es CRO \`a droite)
% ---------------------------------------------------------------
\draw[decorate, decoration={brace, amplitude=5pt, mirror},
      thick, enspblue]
    (4.1, 4.0) -- (4.1, 5.5)
    node[midway, right=6pt, font=\footnotesize, text=enspblue]
        {$\Conf$};
\draw[decorate, decoration={brace, amplitude=5pt, mirror},
      thick, enspgreen]
    (4.1, 2.0) -- (4.1, 3.5)
    node[midway, right=6pt, font=\footnotesize, text=enspgreen]
        {$\Relver$};
\draw[decorate, decoration={brace, amplitude=5pt, mirror},
      thick, navyblue]
    (4.1, 0.0) -- (4.1, 1.5)
    node[midway, right=6pt, font=\footnotesize, text=navyblue]
        {$\Opp$};

\end{tikzpicture}
\caption[Paradigme Gold-Iron-Clay (or, fer, argile)~: architecture \`a trois couches]{%
  Paradigme Gold-Iron-Clay (or, fer, argile)~: architecture \`a trois couches pour les protocoles
  cryptographiques institutionnels.%
}\label{fig:gold_iron_clay}

\smallskip
\begin{minipage}{\linewidth}
  \small\textit{Note :} La couche \textbf{Or} (dor\'ee, $\Conf$) assure le
  masquage cryptographique, la couche \textbf{Fer} (acier, $\Relver$) garantit
  le liage et la v\'erification fiable, tandis que la couche \textbf{Argile} (terre
  cuite, $\Opp$) pr\'esente le contexte sous une forme interpr\'etable~; elle
  n'accro\^it $\Opp$ que si elle transmet une information que l'institution
  ne d\'etient pas d\'ej\`a (Remarque~\ref{rem:statut_sotp}). Les accolades
  droites maintiennent la correspondance avec les propri\'et\'es du Trilemme CRO.
\end{minipage}

\end{figure}
\end{sloppypar}

% -----------------------------------------------------------------------
\subsection{Principes de conception \'equilibr\'ee}
\label{subsec:principes_conception}
% -----------------------------------------------------------------------
\begin{sloppypar}
Un protocole CLO optimal alloue son budget entropique selon les exigences
applicatives. Le budget total d\'efini par
\begin{equation}
H_{\text{total}} = H_{\text{Gold}} + H_{\text{Iron}} + H_{\text{Clay}} \approx 2\lambda + H(\mathcal{C})
\label{eq:entropy_budget_clo}
\end{equation}
vaut par exemple $296$ bits pour $\lambda = 128$ et un contexte de vote
Helios \`a $H(\mathcal{C}) = 40$ bits.

Lorsque l'application priorise la confidentialit\'e ($\Conf_{\min} = 0{,}9$),
l'allocation recommand\'ee affecte $60\%$ du budget \`a la couche Gold
($\approx 180$ bits), $30\%$ \`a Iron ($\approx 90$ bits) et les
$10\%$ restants \`a Clay (compression agressive du contexte). Lorsqu'au
contraire l'opposabilit\'e prime ($\Opp_{\min} = 0{,}8$), les proportions
s'inversent, avec $20\%$ pour Gold (masquage minimal), $30\%$ pour Iron
(binding standard) et $50\%$ pour Clay (pr\'eservation riche du contexte).

L'ind\'ependance modulaire des couches est n\'ecessaire, car Gold, Iron et Clay
doivent pouvoir \^etre substitu\'es sans refonte compl\`ete~\cite{canetti2001,Goldreich02}. Remplacer le PRG
par AES-GCM n'affecte pas Iron ni Clay~; remplacer Pedersen par un
accumulateur RSA laisse Gold et Clay intacts~; remplacer le hachage
contextuel par une compression JL n'impacte pas les deux premi\`eres
couches. Cette modularit\'e facilite l'adaptation aux \'evolutions
cryptographiques (post-quantique, homomorphe) et r\'eglementaires~\cite{NISTPQC,Kyber2022,Dilithium2022}.
\end{sloppypar}

% -----------------------------------------------------------------------
\subsection{Construction~: le Semantic One-Time Pad (S-OTP)}
\label{subsec:construction_sotp}
% -----------------------------------------------------------------------
\begin{sloppypar}
Nous montrons que le paradigme Gold-Iron-Clay est instanciable
dans la classe des protocoles bas\'es XOR, via la construction
\textbf{Semantic One-Time Pad (S-OTP)}, prototype de la famille SOTP
\'etendue \`a la section~\ref{sec:sotp_construction}.

\begin{theorem}[Propri\'et\'es de s\'ecurit\'e S-OTP]\label{thm:proprietes_sotp}
Sous hypoth\`eses PRG (g\'en\'erateur pseudo-al\'eatoire) et difficult\'e du
logarithme discret (DLog) dans $\mathbb{G}$~\cite{BlumMicali84,Yao82}, on a
\begin{align}
\Conf(\text{S-OTP}) &= 1 - \negl(\lambda) \\
\Relver(\text{S-OTP}) &= 1 - \negl(\lambda) \\
\Opp(\text{S-OTP}) &\leq \frac{H(h \mid C)}{H(V|C)}
\end{align}
o\`u $h$ est la composante auxiliaire de la preuve qui porte l'information
contextuelle.
\end{theorem}

\begin{remark}[Port\'ee du Th\'eor\`eme~\ref{thm:proprietes_sotp}]\label{rem:portee_sotp}
La fiabilit\'e \'etablie est la fiabilit\'e de v\'erification $\Relver$~; la
fiabilit\'e de reconstruction v\'erifie $\Relrec \leq (1+\kappa/\lambda)/\lambda$
d\`es que $I(M;\Sigma|C) \leq \kappa/\lambda$ (Remarque~\ref{rem:rel_ver}). Si $h$ est
une fonction du contexte public $C$, alors $H(h|C) = 0$ et $\Opp(\text{S-OTP}) = 0$.
L'opposabilit\'e de S-OTP suppose donc que $h$ transmette une information
contextuelle que l'institution ne d\'etient pas d\'ej\`a (contexte d'esp\`ece
priv\'e). L'expression $H(\mathcal{C})/\lambda$ de~\cite{MINKA2026CLO} correspond \`a la
convention $H(h|C) = H(\mathcal{C})$ et $H(V|C) \approx \lambda$.
\end{remark}

\begin{Preuve}[Th\'eor\`eme~\ref{thm:proprietes_sotp}~: Confidentialit\'e]\label{prv:sotp_conf}

\begin{cmtEtape}{Objectif.}
Montrer que $\Conf(\text{S-OTP}) = 1 - \negl(\lambda)$. Le masquage
one-time pad rend le message chiffr\'e indiscernable, propri\'et\'e dont
h\'erite l'engagement Pedersen.
\end{cmtEtape}

Par la construction du masquage one-time pad, $M^* = m \oplus \mathsf{mask}_m$
o\`u $\mathsf{mask}_m \xleftarrow{\$} \{0,1\}^{|m|}$. Pour tout couple
$(m_0, m_1) \in \mathcal{M}$, on a
\begin{equation}
\TV(P_{M^*|M=m_0}, P_{M^*|M=m_1}) = 0
\end{equation}
\cmtFormule{Le masque uniforme rend les deux distributions identiques.}
\cmt{Un XOR avec un al\'ea uniforme efface toute trace du message.}

car les deux distributions sont uniformes sur $\{0,1\}^{|m|}$.
\cmt{Quel que soit $m$, $m \oplus \mathsf{mask}$ est uniforme.}

L'engagement $c_\Sigma = \mathsf{Com}(M^*)$ h\'erite de cette
indiscernabilit\'e, car l'engagement de Pedersen est \emph{parfaitement} masquant
sans hypoth\`ese calculatoire~\cite{Pedersen91}. Il en r\'esulte que
\begin{equation}
\Conf(\text{S-OTP}) = 1 - \max_{m_0,m_1} \TV(P_{\Sigma|M=m_0}, P_{\Sigma|M=m_1})
= 1 - \negl(\lambda)
\end{equation}
\cmtFormule{La confidentialit\'e atteint son maximum \`a une quantit\'e n\'egligeable pr\`es.}
\cmt{Pedersen cache parfaitement $M^*$, inconditionnellement~; seul le caract\`ere liant repose sur DLog.}
\end{Preuve}

\begin{Preuve}[Th\'eor\`eme~\ref{thm:proprietes_sotp}~: Fiabilit\'e]\label{prv:sotp_rel}

\begin{cmtEtape}{Objectif.}
Montrer que $\Relver(\text{S-OTP}) = 1 - \negl(\lambda)$. Il faut
\'etablir deux propri\'et\'es, la compl\'etude (une preuve honn\^ete est accept\'ee) et
la soundness (une preuve forg\'ee est rejet\'ee).
\end{cmtEtape}

La compl\'etude suit directement de la compl\'etude du sch\'ema
d'engagement. Pour tout $(m, c)$, on a
\begin{equation}
\Pr[\mathsf{Verify}(pk, \mathsf{Prove}(sk,m,c), c) = 1]
= 1 - \delta_{\mathsf{Com}}
\end{equation}
\cmtFormule{Une preuve honn\^ete est accept\'ee sauf erreur du sch\'ema d'engagement.}
o\`u $\delta_{\mathsf{Com}} = \negl(\lambda)$ est l'erreur d'engagement.
\cmt{L'erreur d'engagement est n\'egligeable en $\lambda$.}

Pour \'etablir la \emph{soundness}, supposons qu'un adversaire $\mathcal{A}$ produise
$\sigma^* = (c_\Sigma^*, h^*, \ldots)$ avec $\mathsf{Verify}(\sigma^*, c) = 1$
sans passer par $\mathsf{Prove}$. Alors $\mathcal{A}$ aurait ouvert
l'engagement $c_\Sigma^*$ sur un contenu valide~; cela contredit la
propri\'et\'e de \emph{binding} calculatoire du sch\'ema de Pedersen sous
l'hypoth\`ese DLog (probabilit\'e $\leq \negl(\lambda)$)~\cite{Pedersen91}. Donc
$\Relver(\text{S-OTP}) = 1 - \negl(\lambda)$.
\cmt{Pour forger, il faudrait ouvrir un engagement sur deux valeurs distinctes.}
\cmt{Cela \'equivaut \`a r\'esoudre le logarithme discret, ce qui est difficile.}
\end{Preuve}

\begin{Preuve}[Th\'eor\`eme~\ref{thm:proprietes_sotp}~: Opposabilit\'e]\label{prv:sotp_opp}

\begin{cmtEtape}{Objectif.}
Montrer que $\Opp(\text{S-OTP}) \leq H(h|C)/H(V|C)$.
L'information pertinente pour le verdict ne peut transiter que par la
composante auxiliaire $h$.
\end{cmtEtape}

L'information auxiliaire $h$ encode $\ell_{\text{aux}} = H(\mathcal{C})$
bits de contexte. Par construction, on a
\begin{equation}
I(V;\Sigma|C) \leq I(V;h|C) = H(h|C) - H(h|V,C) \leq H(h|C)
\end{equation}
\cmtFormule{Seule la composante $h$ porte l'information utile au verdict.}
\cmt{Le masque $M^*$ est ind\'ependant des verdicts.}
\cmt{L'engagement $c_\Sigma$ est parfaitement cach\'e.}

En divisant par $H(V|C)$, on obtient
\begin{equation}
\Opp(\text{S-OTP}) = \frac{I(V;\Sigma|C)}{H(V|C)}
\leq \frac{H(h|C)}{H(V|C)}
\end{equation}
\cmtFormule{L'opposabilit\'e est born\'ee par ce que $h$ apprend \`a l'institution au-del\`a du contexte public.}
\cmt{Si $h$ est une fonction de $C$, alors $H(h|C) = 0$ (voir la Remarque~\ref{rem:portee_sotp}).}
\end{Preuve}

\begin{remark}[Optimalit\'e de S-OTP dans la classe OTP~: argument heuristique]\label{thm:optimalite_sotp}
Parmi les protocoles qui masquent le message par un one-time pad et
n'encodent le contexte que dans une composante auxiliaire $h$, S-OTP
semble r\'ealiser le meilleur compromis, puisque tout protocole de ce type
consacre $|m| = \lambda$ bits d'al\'ea au masquage pour la s\'ecurit\'e
parfaite et voit son opposabilit\'e born\'ee par $H(h|C)/H(V|C)$ (m\^eme argument
que ci-dessus). Une preuve d'optimalit\'e demanderait de d\'efinir
formellement la classe OTP et d'y d\'emontrer une borne sup\'erieure
uniforme~; nous la laissons ouverte. Toute construction atteignant une
opposabilit\'e sup\'erieure doit soit (i) abandonner la confidentialit\'e
information-th\'eorique au profit d'hypoth\`eses calculatoires, soit (ii) utiliser
des protocoles interactifs multi-tours, soit (iii) employer des primitives
diff\'erentes, au-del\`a du paradigme OTP.
\end{remark}

\subsection{M\'etriques CLO atteintes}
\label{subsec:metriques_sotp}

Pour les param\`etres typiques $\lambda = 128$, $\Delta = 50$ bits,
$H(\mathcal{C}) = 30$ bits (r\'egime contexte rare), on obtient
\begin{align}
\Conf(\text{S-OTP}) &\approx 0{,}999 \label{eq:sotp_conf_numerique} \\
\Relver(\text{S-OTP}) &\approx 0{,}999 \label{eq:sotp_rel_numerique} \\
\Opp(\text{S-OTP}) &\approx \min\!\left\{1,\; \frac{50}{30}\right\} = 1 \label{eq:sotp_opp_numerique}
\end{align}

Ces valeurs sont des estimations th\'eoriques calcul\'ees sous les hypoth\`eses
PRG et DLog~\cite{Goldreich01,KatzLindell14}, sans mesure empirique
contr\^ol\'ee. Le produit vaut alors
\begin{equation}
\Conf \cdot \Relver \cdot \Opp \approx 0{,}998
\label{eq:sotp_produit}
\end{equation}
Cette valeur porte sur la fiabilit\'e de v\'erification et sur une
opposabilit\'e calcul\'ee sous la convention $\Opp = \min\{1,\Delta/H(\mathcal{C})\}$,
qui suppose que $h$ transmet un contexte d'esp\`ece inconnu de l'institution~;
elle n'est pas comparable \`a la borne du
Th\'eor\`eme~\ref{thm:trilemme_cro_conditionnel}. En fiabilit\'e de
reconstruction, la confidentialit\'e statistique de S-OTP impose
$\Relrec \leq (1+\kappa/\lambda)/\lambda$ (Remarque~\ref{rem:rel_ver}),
tandis que~\cite{MINKA2026CLO} donne un produit de l'ordre de $H(\mathcal{C})/\lambda$.
Ce produit \'elev\'e n'est r\'ealisable qu'en r\'egime contexte rare ($H(\mathcal{C}) \leq \Delta$).
D\`es que le contexte s'enrichit ($H(\mathcal{C}) \gg \Delta$), l'opposabilit\'e
s'effondre \`a $\Opp \approx \Delta / H(\mathcal{C}) \ll 1$, ce qui motive
les variantes SOTP-Base et SOTP-High de la section~\ref{sec:sotp_construction}.

\begin{boiteArbitre}
L'arbitre comprend que $0{,}998$ correspond au meilleur cas. D\`es que le
contexte s'enrichit, l'opposabilit\'e s'effondre et d'autres constructions
deviennent n\'ecessaires.
\end{boiteArbitre}

La remarque~\ref{thm:optimalite_sotp} indique que S-OTP serait optimal
dans sa classe, sans le d\'emontrer. Cette optimalit\'e, m\^eme \'etablie,
resterait restreinte, car les variantes SOTP-Base et SOTP-High utilisent
des techniques au-del\`a des protocoles XOR (compression JL, zk-SNARKs)
et surpassent S-OTP dans les r\'egimes appropri\'es.
\end{sloppypar}

% -----------------------------------------------------------------------